\documentclass[]{article}

\usepackage{amsmath}
\usepackage{multirow}
\usepackage{natbib}
\usepackage{graphicx} 


\begin{document}

In this paper, we addressed the significant computational challenge posed by the need for large ensembles of N-body simulations in modern cosmology. The high computational cost of traditional CPU-based codes creates a bottleneck for analyses requiring robust statistical error estimation. To overcome this, we developed and validated a cosmological Particle-Mesh (PM) N-body simulation accelerated on a Graphics Processing Unit (GPU) using the NVIDIA Warp framework.

We conducted an ensemble of ten independent simulations, each evolving $512^3$ particles in a $(1000 \, \text{Mpc}/h)^3$ volume. The simulations were initialized at $z=127$ using second-order Lagrangian Perturbation Theory (2LPT) and evolved to $z=0$. Our analysis focused on two key aspects: computational performance and physical accuracy. The results demonstrate a profound performance gain, with the wall-clock time for a single realization reduced from hours on a CPU to approximately 20 seconds on a GPU. This enables the generation of large simulation suites in a matter of minutes.

To validate the physical accuracy, we compared the ensemble-averaged matter power spectrum against theoretical predictions. Our GPU-accelerated simulation accurately reproduces the reference power spectrum on large cosmological scales ($k < 0.03 \, \text{h/Mpc}$), with agreement within 5\%. At smaller scales ($k > 0.1 \, \text{h/Mpc}$), the simulation systematically underpredicts power, a well-understood and expected consequence of the force resolution being limited by the grid spacing inherent to the PM method.

We have learned that GPU acceleration with modern frameworks like NVIDIA Warp provides a viable and highly efficient solution for generating cosmological simulations. The validation confirms that our implementation is a powerful tool for applications where large-scale accuracy is paramount, such as studies of the Baryon Acoustic Oscillations or the construction of covariance matrices for galaxy surveys. While not suited for resolving small-scale non-linear structures, this approach effectively removes a major computational barrier, facilitating the rapid production of the large simulation ensembles required by next-generation cosmological analyses.

\end{document}
                